\documentclass[a4paper,11pt]{article}
\usepackage{iftex}
\iftutex
  % If compiling with LuaLaTeX or XeLaTeX
  \usepackage{fontspec}
\else
  % If compiling with pdfLaTeX
  \usepackage[utf8]{inputenc}
  \usepackage[T1]{fontenc}
\fi
\usepackage[english]{babel}
\usepackage{geometry}
\geometry{hmargin=2.5cm, vmargin=2.5cm}
\usepackage{hyperref}
\usepackage{xcolor}
\usepackage{listings}

\hypersetup{
    colorlinks=true,
    linkcolor=blue,
    filecolor=magenta,      
    urlcolor=cyan,
}

\lstset{
  language=[LaTeX]TeX,
  basicstyle=\ttfamily\small,
  keywordstyle=\color{blue},
  commentstyle=\color{gray},
  frame=single,
  backgroundcolor=\color{black!5},
  breaklines=true,
  postbreak=\mbox{\textcolor{gray}{$\hookrightarrow$}\space}
}

\title{\textbf{The \texttt{archery-trainer-sets} Package}}
\author{deogracia}
\newcommand{\pkgversion}{1.1.0-dev}
\date{\today \quad -- \quad Version \pkgversion}

\begin{document}
\maketitle

\begin{abstract}
The \texttt{archery-trainer-sets} package provides a set of \LaTeX{} classes and packages for archery trainers. It aims to standardize and facilitate the creation of documents useful for preparing and following training sessions.
\end{abstract}

\tableofcontents
\vspace{1cm}

\section{Introduction}
This package was designed to meet the specific needs of archery trainers (initiation, progression, competition). Currently, it includes the \texttt{ficheseance} class, which allows generating clear, structured, and print-ready session sheets (Landscape format, 7 columns).

The package is distributed under the \textbf{LPPL 1.3c} license.

\section{The \texttt{ficheseance} class}

\subsection{Loading the class}
The \texttt{ficheseance} class is based on the standard \texttt{article} class. It automatically sets the document to A4 size, landscape orientation, with a 10pt font.
To use it, simply declare:
\begin{lstlisting}
\documentclass{ficheseance}
\end{lstlisting}

The class natively supports French and English via the \texttt{translations} package. Just load \texttt{babel} or \texttt{polyglossia} with your desired language, and the headers will translate automatically.

\subsection{Defining metadata}
Before the \verb|\begin{document}| or right after, you must define the information of your session using the following commands:
\begin{itemize}
    \item \verb|\numfiche{...}|: Number of the sheet or session.
    \item \verb|\datefiche{...}|: Date of the session.
    \item \verb|\archer{...}|: Name of the archer or group.
    \item \verb|\theme{...}|: General theme of the session.
    \item \verb|\objectifseance{...}|: Main objective of the session.
\end{itemize}

\subsection{Header generation}
Once the information is defined, use the \verb|\makeficheheader| command to display the summary box at the top of the first page.

\subsection{The session table}
The core of the sheet is the \texttt{tableauSeance} environment. It deploys a multi-page table (thanks to \texttt{xltabular}) consisting of 7 columns specific to sports training:
\begin{enumerate}
    \item \textbf{Goal}
    \item \textbf{Time}
    \item \textbf{Situations}
    \item \textbf{Tools}
    \item \textbf{Instructions}
    \item \textbf{Success criteria}
    \item \textbf{Evaluations}
\end{enumerate}

To insert a row (a situation or exercise) into the table, use the \verb|\ligneSeance| command, which takes exactly 7 arguments:
\begin{lstlisting}
\begin{tableauSeance}
    \ligneSeance
        {Warm-up} % Goal
        {15 min} % Time
        {Joint and muscle awakening} % Situation
        {Elastic bands} % Tools
        {Slow and controlled movements} % Instructions
        {No pain} % Criteria
        {OK} % Evaluation
\end{tableauSeance}
\end{lstlisting}

\section{Support and Contribution}
The source code and full examples are available on the GitHub repository:\\ \url{https://github.com/deogracia/archery-trainer-sets}.\\
Contributions as \textit{Pull Requests} are welcome.

\end{document}